\documentclass{aastex631}

\usepackage{amsmath,amssymb}
\usepackage{natbib}

\newcommand{\dd}{\mathrm{d}}
\newcommand{\pp}{\partial}
\newcommand{\kbbar}{k_{\rm B}/\mathrm{amu}}
\newcommand{\Cp}{C_P}
\newcommand{\Cv}{C_V}

\shorttitle{ORCHARD Hydrostatic Methods}
\shortauthors{Tejada Arevalo}

\begin{document}

\title{Henyey Relaxation for Hydrostatic Equilibrium in the ORCHARD Planetary Evolution Code}

\author{Roberto Tejada Arevalo}

\begin{abstract}
We describe the Henyey relaxation method used by the ORCHARD planetary
evolution code to solve for the hydrostatic structure of planetary
interiors.  The solver determines pressure $P(m)$ and radius $r(m)$ on a
Lagrangian mass grid of $N$ shells by simultaneously satisfying
hydrostatic equilibrium (including centrifugal support from rigid-body
rotation) and mass continuity.  The $2N$ coupled nonlinear equations are
cast in logarithmic variables $(\ln P,\ln r)$ and solved with an
under-damped Newton--Raphson iteration whose Jacobian has a sparse,
block-tridiagonal structure.  The equation of state is evaluated
separately for each compositional layer---hydrogen--helium envelope,
silicate mantle, and iron core---with entropy, helium fraction, and
metallicity held fixed during the structural relaxation.  We also
describe the boundary conditions, rotational corrections, moment-of-%
inertia feedback, and the optional Theory-of-Figures treatment of
gravitational harmonics.
\end{abstract}

\keywords{planets and satellites: interiors --- planets and satellites:
physical evolution --- methods: numerical}

% ======================================================================
\section{Introduction}\label{sec:intro}
% ======================================================================

Stellar and planetary structure calculations have long relied on the
Henyey relaxation scheme \citep{Henyey1964}, which solves the full set
of structure equations simultaneously on a one-dimensional grid rather
than integrating them as a boundary-value shooting problem.  The global
Newton--Raphson iteration linearizes the residual of the discretized
equations and inverts the resulting block-tridiagonal Jacobian, achieving
quadratic convergence near the solution.

In ORCHARD, the hydrostatic equilibrium solver accepts thermodynamic
profiles---entropy $S(m)$, helium mass fraction $Y(m)$, and heavy-%
element mass fraction $Z(m)$---from the transport module and returns
self-consistent pressure, radius, density, and temperature profiles that
satisfy mechanical equilibrium.  These updated profiles are then returned
to the transport solver, closing the structure--transport coupling loop
\citep{Sur2024}.

The planet is divided into three structural regions indexed by mass
coordinate $m$:
\begin{enumerate}
  \item \textbf{Envelope} (cells $0$ to $k_{\rm core}-1$): hydrogen--helium
    mixture with dissolved heavy elements.
  \item \textbf{Rocky mantle} (cells $k_{\rm core}$ to $k_{\rm core,Fe}-1$):
    silicate material (e.g.\ MgSiO$_3$, Mg$_2$SiO$_4$).
  \item \textbf{Iron core} (cells $k_{\rm core,Fe}$ to $N-1$): metallic
    iron or iron alloy.
\end{enumerate}
Special cases include coreless gas giants ($k_{\rm core}=N$), planets
with no iron core ($k_{\rm core,Fe}=N$), and bare-rock super-Earths
($k_{\rm core}=0$).

% ======================================================================
\section{State Vector}\label{sec:state}
% ======================================================================

The solver works in logarithmic variables.  For $N$ mass shells the
state vector $\mathbf{Y}$ has $2N$ components, interleaved as
\begin{equation}\label{eq:statevec}
  \mathbf{Y}
  = \bigl(\ln P_0,\;\ln r_0,\;\ln P_1,\;\ln r_1,\;\ldots,\;
          \ln P_{N-1},\;\ln r_{N-1}\bigr)^{\!\top}\!,
\end{equation}
where $P_k$ is the pressure at the center of zone~$k$ and $r_k$ is the
radius of its outer boundary.  Even-indexed elements carry
$\ln P$ and odd-indexed elements carry $\ln r$.

The thermodynamic quantities that close the system---temperature
$T_k(S_k, P_k, Y_k, Z_k)$ and density $\rho_k(S_k, P_k, Y_k,
Z_k)$---are evaluated from the equation of state (EOS) at each
iteration but are \emph{not} part of the state vector: they are treated
as auxiliary variables updated from the current $(S, P)$.

% ======================================================================
\section{Structure Equations}\label{sec:equations}
% ======================================================================

The residual vector $\mathbf{A}$ contains $2N$ equations---one pressure
(hydrostatic equilibrium) equation and one radius (mass continuity)
equation per zone.  We define auxiliary logarithmic variables
\begin{equation}
  x_k \equiv \ln r_k, \qquad
  y_k \equiv \ln P_k, \qquad
  q_k \equiv \ln \rho_k.
\end{equation}
The mass grid boundaries $m_k$ ($k = 0,\ldots,N$) are fixed input with
$m_0 = M_{\rm planet}$ at the surface and $m_N = 0$ at the center.
The cell-center mass is $\bar{m}_k = \tfrac{1}{2}(m_k + m_{k+1})$.

\subsection{Hydrostatic Equilibrium}\label{sec:hydrostatic}

The continuous hydrostatic equilibrium equation for a rotating,
self-gravitating sphere (in the spherically averaged approximation) is
\begin{equation}\label{eq:hse_cont}
  \frac{\dd P}{\dd m}
  = -\frac{G\,m}{4\pi r^4}
    + \frac{\omega(t)^2\, r}{6\pi},
\end{equation}
where $\omega(t)$ is the time-dependent angular velocity of rigid-body
rotation, determined at each timestep by angular momentum conservation
(Section~\ref{sec:ang_mom}), and the second term is the volume-averaged
centrifugal acceleration of a uniformly rotating sphere.  Equation~\eqref{eq:hse_cont} is discretized
differently for the top zone, interior zones, and inner boundary.

\subsubsection{Outer Boundary ($k=0$)}

The outermost zone is matched to an imposed atmospheric pressure
$P_{\rm atm}$ (typically 1~bar):
\begin{align}\label{eq:A0}
  A_0 &= -P_{\rm atm} + P_0
        - \frac{G}{4\pi}\,m_0\,\frac{\Delta m_0}{2}\,e^{-4x_0}
        \notag\\
      &\quad + \frac{\Delta m_0}{2}\,\frac{\omega^2}{6\pi}\,e^{-x_0},
\end{align}
where $\Delta m_0 \equiv m_0 - m_1$ is the mass of the outermost shell.
At convergence $A_0 = 0$ enforces $P_0 = P_{\rm atm} +
\Delta P_{\rm grav} - \Delta P_{\rm cent}$.

\subsubsection{Interior Zones ($k = 1,\ldots,N-2$)}

For each interior zone the pressure difference between adjacent shells
balances gravity and centrifugal force:
\begin{align}\label{eq:Ak_interior}
  A_{2k} &= y_{k-1} - y_k
           + \frac{G}{4\pi}\,\frac{\Delta m_k^{\star}}{2}\,m_k\,
             e^{-y_k - 4x_k}
           \notag\\
         &\quad
           - \frac{\Delta m_k^{\star}}{2}\,\frac{\omega_k^2}{6\pi}\,
             e^{-x_k - \frac{1}{2}(y_{k-1}+y_k)},
\end{align}
where $\Delta m_k^{\star} \equiv m_{k-1} - m_{k+1}$ spans two shell
boundaries.

The exponential factors arise from writing the gravitational
acceleration in logarithmic variables:
\begin{equation}
  \frac{G\,\bar{m}_k}{4\pi\,r_k^4}
  = \frac{G\,\bar{m}_k}{4\pi}\,e^{-4\ln r_k}
  = \frac{G\,\bar{m}_k}{4\pi}\,e^{-4\,x_k}.
\end{equation}
The centrifugal term
$e^{-x_k - \frac{1}{2}(y_{k-1}+y_k)}$ provides a weak pressure--%
dependent correction through the geometric-mean pressure factor; this
coupling appears naturally in the finite-difference approximation of the
product $\omega^2\,r/P$.

\subsubsection{Inner Boundary ($k=N-1$)}

At the center the discretization switches to a form that enforces
regularity.  The innermost pressure residual is
\begin{equation}\label{eq:AN}
  A_{2(N-1)} = P_{N-2} - P_{N-1}
              + \frac{G}{2}\left(\frac{4\pi}{3}\right)^{\!1/3}
                m_{N-1}^{2/3}\,\rho_{N-1}^{4/3},
\end{equation}
which approximates the central pressure increment using a uniform-%
density sphere of mass $m_{N-1}$.  Here $P_{N-2}$ and $P_{N-1}$ are
the exponential pressures (not their logarithms); this avoids
logarithmic singularities at the center.

\subsection{Mass Continuity}\label{sec:mass_cont}

The mass continuity equation in Lagrangian form is
\begin{equation}\label{eq:mass_cont}
  \frac{\dd r}{\dd m} = \frac{1}{4\pi r^2 \rho}.
\end{equation}
This is discretized as follows.

\subsubsection{Non-Central Zones ($k = 0,\ldots,N-2$)}

\begin{equation}\label{eq:radius_eq}
  A_{2k+1} = x_k - x_{k+1}
            - \frac{m_k - m_{k+1}}{4\pi}\,
              e^{-\frac{3}{2}(x_k + x_{k+1}) - q_k},
\end{equation}
which, upon setting $A_{2k+1}=0$, gives the shell thickness
$\Delta\ln r$ consistent with the enclosed mass.  The exponent
$-\tfrac{3}{2}(x_k + x_{k+1})$ represents $(r_k\,r_{k+1})^{-3/2}$, a
geometric-mean area factor that symmetrizes the scheme.

\subsubsection{Central Zone ($k=N-1$)}

The innermost zone radius is fixed by the enclosed mass and
the local density:
\begin{equation}\label{eq:radius_center}
  A_{2N-1} = x_{N-1}
            - \frac{1}{3}\!\left[\ln\!\left(\frac{3\,m_{N-1}}{4\pi}\right)
              - q_{N-1}\right],
\end{equation}
which enforces $r_{N-1} = (3\,m_{N-1} / 4\pi\rho_{N-1})^{1/3}$, the
radius of a uniform-density sphere containing the innermost mass shell.

% ======================================================================
\section{Jacobian Matrix}\label{sec:jacobian}
% ======================================================================

The Newton correction $\delta\mathbf{Y}$ at each iteration satisfies
\begin{equation}\label{eq:newton}
  \mathbf{J}\,\delta\mathbf{Y} = \mathbf{A},
\end{equation}
where $\mathbf{J} = \pp\mathbf{A}/\pp\mathbf{Y}$ is the $2N\times 2N$
Jacobian matrix.  Because each residual equation couples at most three
adjacent zones, $\mathbf{J}$ has a sparse, block-tridiagonal structure
with $2\times 2$ blocks.

\subsection{Interior-Zone Derivatives}

For an interior zone $k\in\{1,\ldots,N-2\}$, the non-zero Jacobian
entries of the pressure equation ($i = 2k$) are:
\begin{align}
  J_{i,\,i-2}
  &= \frac{\pp A_{2k}}{\pp y_{k-1}}
   = 1 + \frac{\omega_k^2\,\Delta m_k^{\star}}{12\pi}\,
     e^{-x_k - \frac{1}{2}(y_{k-1}+y_k)},
  \label{eq:J_pk_left}
\\[4pt]
  J_{i,\,i}
  &= \frac{\pp A_{2k}}{\pp y_k}
   = -1 - \frac{G\,\Delta m_k^{\star}\,m_k}{8\pi}\,
     e^{-y_k - 4x_k}
  \notag\\
  &\quad + \frac{\omega_k^2\,\Delta m_k^{\star}}{12\pi}\,
     e^{-x_k - \frac{1}{2}(y_{k-1}+y_k)},
  \label{eq:J_pk_self}
\\[4pt]
  J_{i,\,i+1}
  &= \frac{\pp A_{2k}}{\pp x_k}
   = -\frac{G\,\Delta m_k^{\star}\,m_k}{\pi}\,
     e^{-y_k - 4x_k}
  \notag\\
  &\quad + \frac{\omega_k^2\,\Delta m_k^{\star}}{6\pi}\,
     e^{-x_k - \frac{1}{2}(y_{k-1}+y_k)}.
  \label{eq:J_pk_right}
\end{align}
Here $J_{i,\,i-2}$ couples to $\ln P_{k-1}$ (two state-vector
positions to the left), $J_{i,\,i}$ is the diagonal ($\ln P_k$), and
$J_{i,\,i+1}$ couples to $\ln r_k$ (one position to the right).

The radius equation ($i = 2k+1$) has two non-zero entries:
\begin{align}
  J_{i,\,i}
  &= \frac{\pp A_{2k+1}}{\pp x_k}
   = 1 + \frac{3(m_k - m_{k+1})}{8\pi}\,
     e^{-\frac{3}{2}(x_k+x_{k+1}) - q_k},
  \label{eq:J_rk_self}
\\[4pt]
  J_{i,\,i+2}
  &= \frac{\pp A_{2k+1}}{\pp x_{k+1}}
   = -1 + \frac{3(m_k - m_{k+1})}{8\pi}\,
     e^{-\frac{3}{2}(x_k+x_{k+1}) - q_k}.
  \label{eq:J_rk_right}
\end{align}

\subsection{Outer Boundary Derivatives}

At the surface ($k=0$) the Jacobian of the pressure equation has
contributions from $\ln P_0$ and $\ln r_0$:
\begin{align}
  J_{0,0} &= P_0, \label{eq:J_surf_P} \\
  J_{0,1} &= \frac{G\,m_0\,\Delta m_0}{\pi}\,e^{-4x_0}
             - \frac{\Delta m_0\,\omega_0^2}{12\pi}\,e^{-x_0}.
  \label{eq:J_surf_r}
\end{align}
The radius equation at $k=0$ takes the same form as the interior
(Equations~\ref{eq:J_rk_self}--\ref{eq:J_rk_right} with $k=0$).

\subsection{Inner Boundary Derivatives}

At the center ($k=N-1$):
\begin{equation}
  J_{2(N-1),\,2(N-2)} = P_{N-2},\qquad
  J_{2(N-1),\,2(N-1)} = -P_{N-1},
\end{equation}
reflecting the exponential pressure terms in
Equation~\eqref{eq:AN}, and
\begin{equation}
  J_{2N-1,\,2N-1} = 1,
\end{equation}
since the central radius equation~\eqref{eq:radius_center} depends
only on $x_{N-1}$, with $q_{N-1}$ treated as a known quantity at
each iteration.

\subsection{Sparse Storage}

The full Jacobian is stored in compressed sparse column (CSC) format
using SciPy's \texttt{csc\_array}.  All non-zero entries are computed
vectorially over the zone index~$k$, yielding an assembly cost that
scales linearly in~$N$.  The sparse direct solve (LU factorization via
\texttt{spsolve}) exploits the narrow bandwidth of the matrix.

% ======================================================================
\section{Iteration Scheme}\label{sec:iteration}
% ======================================================================

\subsection{Damped Newton--Raphson}

Given an initial guess $\mathbf{Y}^{(0)}$ (typically from the previous
timestep), the iteration proceeds as follows:
\begin{enumerate}
  \item \textbf{EOS update.}  Evaluate $T_k$ and $\rho_k$ from the
    current $(S_k, P_k, Y_k, Z_k)$ using the appropriate EOS for each
    compositional layer (Section~\ref{sec:eos}).
  \item \textbf{Residual and Jacobian.}  Compute $\mathbf{A}$ and
    $\mathbf{J}$ from the current $(x, y, q)$.
  \item \textbf{Linear solve.}  Solve $\mathbf{J}\,\delta\mathbf{Y}
    = \mathbf{A}$ using sparse LU decomposition.
  \item \textbf{Damped update.}  Apply the correction with a fixed
    damping factor $\alpha$:
    \begin{equation}\label{eq:update}
      \mathbf{Y}^{(n+1)} = \mathbf{Y}^{(n)} - \alpha\,\delta\mathbf{Y},
      \qquad \alpha = 0.4.
    \end{equation}
  \item \textbf{Convergence check.}  Evaluate the relative error
    \begin{equation}\label{eq:convergence}
      \varepsilon
      = \max_{k}\,
        \frac{\bigl|Y_k^{(n+1)} - Y_k^{(n)}\bigr|}
             {\bigl|Y_k^{(n)}\bigr| + 10^{-30}},
    \end{equation}
    and declare convergence when $\varepsilon < \varepsilon_{\rm tol}$
    (default $10^{-10}$).
\end{enumerate}
The under-relaxation factor $\alpha = 0.4$ prevents overshoot in the
early iterations when the initial guess may be far from the solution;
for near-converged states the effective convergence rate remains
rapid.  A maximum of 2000 iterations is permitted before the step
is rejected.

\subsection{Failure Handling}

Two guards protect against singular or divergent iterations.  If the
sparse solve yields non-finite values in $\delta\mathbf{Y}$ (indicating
a singular Jacobian), or if $\mathbf{Y}^{(n+1)}$ itself contains NaN,
the solver immediately writes NaN to the pressure and radius arrays and
exits.  The calling evolution loop detects the NaN sentinel and rejects
the timestep, reverting to a smaller $\Delta t$ before re-attempting.

% ======================================================================
\section{Equation of State Interface}\label{sec:eos}
% ======================================================================

The thermodynamic closure---converting $(S, P)$ into $(T, \rho)$---is
performed by three distinct EOS backends, one for each compositional
layer.  Units are CGS internally (pressure in dyn\,cm$^{-2}$,
temperature in~K, density in g\,cm$^{-3}$); conversions to the native
units of each EOS backend are handled at the interface.

\subsection{Envelope (Hydrogen--Helium--Metals)}

Temperature and density for cells $k < k_{\rm core}$ are obtained from
the \texttt{mixtures\_eos} module, which blends the H--He tables of
\citet{Becker2018} and \citet{French2019} with a heavy-element
component at mass fraction~$Z$:
\begin{align}
  \log_{10} T_k &= \texttt{get\_logt\_sp}
    \bigl(S_k,\;\log_{10}P_k,\;Y_k,\;Z_k,\;f_{\rm rock,k}\bigr),
  \label{eq:eos_env_T}
\\
  \log_{10}\rho_k &= \texttt{get\_logrho\_pt}
    \bigl(\log_{10}P_k,\;\log_{10}T_k,\;Y_k,\;Z_k,\;f_{\rm rock,k}\bigr).
  \label{eq:eos_env_rho}
\end{align}
When the tabulated EOS mode is active (\texttt{eos\_tab = True}), the
calls reduce to interpolation on pre-computed
$(S, \log P, Y, Z) \to \log T$ grids, avoiding costly on-the-fly
Gibbs minimization.

\subsection{Rocky Mantle}

For the silicate mantle ($k_{\rm core} \le k < k_{\rm core,Fe}$), the
\texttt{mantle\_eos} module provides
\begin{align}
  T_k &= T(S_k,\;P_k), \label{eq:eos_mantle_T} \\
  \rho_k &= \rho(S_k,\;P_k), \label{eq:eos_mantle_rho}
\end{align}
where pressures are converted from dyn\,cm$^{-2}$ to GPa.  The mantle
EOS is typically a Helmholtz free-energy model for MgSiO$_3$ bridgmanite
or post-perovskite.

\subsection{Iron Core}

For the iron core ($k \ge k_{\rm core,Fe}$), the \texttt{core\_eos}
module provides temperature from the entropy--pressure relation and
density from the pressure--temperature relation:
\begin{align}
  T_k &= T(S_k,\;P_k), \label{eq:eos_fe_T}\\
  \rho_k &= \rho(P_k,\;T_k). \label{eq:eos_fe_rho}
\end{align}
Pressures are converted to Pa for the iron EOS tables, and the
returned density (in kg\,m$^{-3}$) is converted to g\,cm$^{-3}$.

\subsection{NaN Guard for EOS Extrapolation}\label{sec:nan_guard}

When cells lie outside the EOS table domain (e.g.\ low-pressure
surface shells), the returned $T$ or $\rho$ may be NaN.  A
nearest-neighbor fill algorithm replaces non-finite values:
each bad cell is first filled from the nearest inward finite
neighbor; if none exists, the outward direction is searched as a
fallback.  This prevents a single NaN from poisoning the entire
sparse solve.

% ======================================================================
\section{Core Entropy Initialization}\label{sec:core_init}
% ======================================================================

The compact core (mantle + iron) requires separate entropy
initialization since its composition is discontinuous with the
envelope.  ORCHARD supports several strategies.

\subsection{Isothermal Compact-Core Seeding}

When the \texttt{isothermal\_compact\_core} flag is set and both an
envelope and a mantle are present, the entire compact core is seeded
at the temperature of the envelope base:
\begin{equation}\label{eq:iso_seed}
  T_k = T_{k_{\rm core}-1}
  \quad\text{for all}\; k \ge k_{\rm core}.
\end{equation}
The corresponding entropies are then derived from the temperature
and local pressure via the mantle or iron EOS:
\begin{align}
  S_k &= S_{\rm mantle}(P_k,\;T_{k_{\rm core}-1}),
  &k_{\rm core} \le k < k_{\rm core,Fe},
  \label{eq:iso_seed_mantle}\\
  S_k &= S_{\rm Fe}(P_k,\;T_{k_{\rm core}-1}),
  &k \ge k_{\rm core,Fe}.
  \label{eq:iso_seed_fe}
\end{align}

\subsection{Temperature-Continuity Seeding}

When isothermal seeding is not used and the \texttt{env\_s\_start} flag
is set, the mantle is seeded by enforcing temperature continuity at the
envelope--mantle boundary:
\begin{equation}
  T_{k_{\rm core}} = T_{k_{\rm core}-1},
\end{equation}
and then computing a uniform mantle entropy
$S_{\rm mantle} = S(P_{k_{\rm core}}, T_{k_{\rm core}})$ from the
mantle EOS.  Similarly, the iron core inherits the temperature at the
mantle--iron interface:
\begin{equation}
  T_{k_{\rm core,Fe}} = T_{k_{\rm core,Fe}-1},
\end{equation}
with an optional entropy offset $\Delta S_{\rm Fe}$ added to the
iron-core entropy to account for the Clausius--Clapeyron jump at
the phase boundary.

\subsection{Fixed Entropy Seeding}

If neither of the above is activated, the mantle and iron core
are initialized with user-specified uniform entropies
$S_{\rm mantle,0}$ and $S_{\rm Fe,0}$ read from the configuration
file.

% ======================================================================
\section{Rotation and Gravitational Moments}\label{sec:rotation}
% ======================================================================

\subsection{Angular Momentum Conservation}\label{sec:ang_mom}

ORCHARD assumes rigid-body rotation throughout the planet, so that every
mass shell shares a single angular velocity $\omega(t)$ at any given
time.  The total angular momentum is
\begin{equation}\label{eq:Ltot}
  L = I_{\rm struct}\,\omega,
\end{equation}
where $I_{\rm struct}$ is the structural moment of inertia of the
planet.  In the absence of external torques (tidal dissipation, magnetic
braking, mass loss, etc.), angular momentum is conserved:
\begin{equation}\label{eq:Lconserve}
  L = I_{\rm struct}(t)\;\omega(t) = \text{const}.
\end{equation}
As the planet cools and contracts, the density profile $\rho(r)$
evolves, changing $I_{\rm struct}$.  Conservation of $L$ then requires
the angular velocity to adjust at every timestep so that the product
$I_{\rm struct}\,\omega$ remains constant.

Concretely, at initialization the code defines a \emph{reference}
angular momentum from a user-specified rotation period $\mathcal{P}$
and a reference moment of inertia:
\begin{equation}\label{eq:Lref}
  L_{\rm ref} = I_{\rm ref}\,\omega_{\rm ref},
  \qquad
  \omega_{\rm ref} = \frac{2\pi}{\mathcal{P}},
\end{equation}
where
\begin{equation}\label{eq:Iref}
  I_{\rm ref} = C_{\rm MoI}\,M_{\rm planet}\,R_{\rm ref}^2
\end{equation}
is the reference moment of inertia.  Here $C_{\rm MoI}$ is a
dimensionless moment-of-inertia coefficient (e.g.\ $C_{\rm MoI}
\approx 0.254$ for Jupiter), $M_{\rm planet}$ is the total planetary
mass, and $R_{\rm ref}$ is the equatorial reference radius (e.g.\
$R_{\rm Jup}$ or $R_{\rm Sat}$).  This reference state encodes the
\emph{known} present-day angular momentum of the planet.

\subsection{Structural Moment of Inertia}\label{sec:Istruct}

At each Henyey iteration the structural moment of inertia is computed
from the current density and radius profiles.  In the default
(spherically symmetric) approximation,
\begin{equation}\label{eq:Istruct}
  I_{\rm struct} = \frac{8\pi}{3}\int_0^{R}\!r^4\,\rho(r)\,\dd r,
\end{equation}
where $R = r_0$ is the planetary surface radius, $r$ is the radial
coordinate, and $\rho(r)$ is the mass density at radius $r$.  The
integral is evaluated numerically by trapezoidal quadrature over the
$N$ mass shells.  Because $I_{\rm struct}$ depends on the radial
profile, it changes each time the structure is updated, coupling
the rotation rate to the structural relaxation.

\subsection{Time-Dependent Angular Velocity}\label{sec:omega_time}

Given the conserved angular momentum $L_{\rm ref}$ and the current
structural moment of inertia, the angular velocity at each timestep is
\begin{equation}\label{eq:omega}
  \omega(t)
  = \omega_{\rm ref}\,
    \min\!\left(\frac{I_{\rm ref}}{I_{\rm struct}(t)},\;1\right)
  = \frac{2\pi}{\mathcal{P}}\,
    \min\!\left(\frac{I_{\rm ref}}{I_{\rm struct}(t)},\;1\right).
\end{equation}
The $\min(\cdot,\,1)$ clamp ensures that $\omega$ never exceeds the
reference rotation rate $\omega_{\rm ref}$; this prevents unphysically
fast rotation when the structural moment of inertia exceeds the
reference value (for instance, during early hot, inflated stages).

The physical content of Equation~\eqref{eq:omega} is straightforward:
\begin{itemize}
  \item If $I_{\rm struct} > I_{\rm ref}$, the planet is more extended
    than the reference state and $\omega = \omega_{\rm ref}$ (clamped).
  \item If $I_{\rm struct} < I_{\rm ref}$, the planet is more centrally
    concentrated than the reference state.  Angular momentum
    conservation requires $\omega > \omega_{\rm ref}$, but the code
    enforces $\omega = \omega_{\rm ref}\,I_{\rm ref}/I_{\rm struct}$,
    i.e.\ the planet spins up as it contracts.
  \item As $t\to\infty$ and the planet asymptotically approaches its
    present-day structure, $I_{\rm struct} \to I_{\rm ref}$ and
    $\omega \to \omega_{\rm ref}$, recovering the observed rotation
    period.
\end{itemize}

Within a single Henyey relaxation (at fixed time $t$), the angular
velocity is recomputed at every iteration from the latest density
profile.  This provides self-consistent coupling between the
centrifugal support term in the hydrostatic equilibrium equation
(Equation~\ref{eq:hse_cont}) and the evolving structure.

\subsection{Centrifugal Term in Hydrostatic Equilibrium}

The angular velocity enters the hydrostatic balance through the
centrifugal acceleration.  For rigid-body rotation, the
volume-averaged centrifugal acceleration of a uniformly rotating
sphere contributes to Equation~\eqref{eq:hse_cont} as
\begin{equation}\label{eq:cent_accel}
  a_{\rm cent} = \frac{\omega^2\,r}{3},
\end{equation}
where $r$ is the spherical radius.  The factor of $1/3$ arises from
averaging the centrifugal acceleration $\omega^2 r\sin^2\theta$ over
solid angle.  In the discretized pressure equation, this produces the
terms proportional to $\omega^2/(6\pi)$ appearing in
Equations~\eqref{eq:A0} and~\eqref{eq:Ak_interior}.  The angular
velocity $\omega_k$ used in each zone is the global rotation rate from
Equation~\eqref{eq:omega}; the subscript~$k$ emphasizes that the
centrifugal correction is evaluated at the radius of zone~$k$.

\subsection{Theory of Figures}\label{sec:tof}

When the \texttt{I\_correction} flag is enabled, the moment of inertia
and the gravitational harmonics $J_{2n}$ ($n = 1,2,3,4$) are computed
from the fourth-order Theory of Figures \citep[ToF4;][]{Nettelmann2017},
rather than from the approximate integral
Equation~\eqref{eq:Istruct}.  The ToF4 procedure introduces a
smallness parameter
\begin{equation}\label{eq:smallq}
  q = \frac{\omega^2}{\omega_0^2},
  \qquad
  \omega_0 = \sqrt{\frac{G\,M_{\rm planet}}{R^3}},
\end{equation}
where $\omega_0$ is the dynamical frequency, $G$ is the gravitational
constant, and $R$ is the mean planetary radius.  The parameter $q$
measures the ratio of centrifugal to gravitational acceleration at the
equator.

The ToF4 method then solves for the figure functions $s_{2n}(l)$ ($n =
1,2,3,4$) on the mean radius coordinate~$l$ by Jacobi iteration of the
Clairaut--Radau--Darwin integral equations.  These figure functions
describe the departure of level surfaces from spheres.  The
outputs---the gravitational moments $J_2$, $J_4$, $J_6$, $J_8$,
the equatorial and polar radii, and a corrected $I_{\rm struct}$---are
fed back into the Henyey solver through the angular velocity via
Equation~\eqref{eq:omega}.

At very early times ($t < 10^{-3}$~Gyr), the structure is too far from
equilibrium for ToF4 to converge, so the code falls back to the
trapezoidal approximation (Equation~\ref{eq:Istruct}) regardless of the
\texttt{I\_correction} setting.

\subsection{Evolution Loop Coupling}

To summarize the rotation--structure coupling at the level of the
evolution loop:
\begin{enumerate}
  \item At timestep $t^n$, the transport module updates the entropy and
    composition profiles.
  \item The Henyey solver receives these profiles along with the
    angular velocity from the previous timestep, $\omega^{n-1}$.
  \item Within each Henyey iteration $j$:
  \begin{enumerate}
    \item The density profile $\rho_k^{(j)}$ is obtained from the EOS.
    \item $I_{\rm struct}^{(j)}$ is computed from
      Equation~\eqref{eq:Istruct} (or from ToF4 if enabled).
    \item $\omega^{(j)}$ is updated from Equation~\eqref{eq:omega}.
    \item The centrifugal terms in the residual and Jacobian are
      evaluated with $\omega^{(j)}$.
  \end{enumerate}
  \item Upon convergence, the final $\omega^n$ is stored and passed
    to the next timestep.
\end{enumerate}
The angular velocity $\omega$ is stored in the HDF5 output file at
every timestep, enabling post-processing of the rotational evolution
history.

% ======================================================================
\section{Mass Grid}\label{sec:grid}
% ======================================================================

The $N+1$ mass boundaries $m_k$ are constructed by the initialization
module with non-uniform spacing that concentrates zones near the surface
and near compositional boundaries (core--envelope interface).  The
surface zone width is controlled by a parameter $\Delta m_{\rm surf}$
that sets the fractional mass of the outermost shells.  Typical values
are $N = 700$ zones with $\Delta m_{\rm surf}/M_{\rm planet} \sim
10^{-2}$.

The mass grid is fixed throughout the evolution; all structural
variables ($P$, $r$, $\rho$, $T$, $S$) float on the Lagrangian
grid as functions of time.

% ======================================================================
\section{Algorithm Summary}\label{sec:algorithm}
% ======================================================================

The complete Henyey relaxation algorithm for a single call from the
evolution loop is:

\begin{enumerate}
  \item \textbf{Input}: entropy $S_k$, composition $(Y_k, Z_k,
    f_{\rm rock,k})$, mass grid $m_k$, previous structure
    $(\ln P_k^{(0)}, \ln r_k^{(0)})$, previous rotation rate
    $\omega^{(0)}$.

  \item \textbf{Initialize state vector}
    $\mathbf{Y}^{(0)} = (\ln P_0,\,\ln r_0,\,\ldots)$.

  \item \textbf{Iterate} ($n = 0,1,2,\ldots$):
  \begin{enumerate}
    \item Evaluate $T_k^{(n)}$, $\rho_k^{(n)}$ from EOS
      (Section~\ref{sec:eos}).
    \item Seed or update core entropies if $n = 0$ and
      $\texttt{init} = \texttt{True}$
      (Section~\ref{sec:core_init}).
    \item Compute $\omega^{(n)}$ from $I_{\rm struct}$
      (Section~\ref{sec:rotation}).
    \item Assemble residual $\mathbf{A}^{(n)}$ and Jacobian
      $\mathbf{J}^{(n)}$ (Sections~\ref{sec:equations}--\ref{sec:jacobian}).
    \item Solve $\mathbf{J}^{(n)}\,\delta\mathbf{Y}^{(n)} =
      \mathbf{A}^{(n)}$.
    \item Update: $\mathbf{Y}^{(n+1)} = \mathbf{Y}^{(n)} -
      0.4\,\delta\mathbf{Y}^{(n)}$.
    \item Check convergence (Equation~\ref{eq:convergence}).
  \end{enumerate}

  \item \textbf{Final EOS evaluation} to obtain converged $T$, $\rho$.

  \item \textbf{Output}: $r_k$, $P_k$, $\rho_k$, $T_k$, $S_k$,
    surface gravity $g = G\,M_{\rm planet}/r_0^2$,
    angular velocity~$\omega$, gravitational moments $J_{2n}$,
    and shape coefficients.
\end{enumerate}

% ======================================================================
\section{Numerical Properties}\label{sec:properties}
% ======================================================================

\subsection{Scaling}

The per-iteration cost is dominated by the sparse LU solve, which for a
bandwidth-$b$ matrix of size $2N$ scales as $\mathcal{O}(N\,b^2)$.
Since $b = 3$ (the stencil width), the cost is effectively
$\mathcal{O}(N)$.  With $N = 700$ and typical convergence in 10--50
iterations, a single hydrostatic solve requires $\lesssim$\,10\,ms on
modern hardware.

\subsection{Convergence Properties}

For well-initialized problems (structure from the previous timestep),
convergence to $\varepsilon < 10^{-10}$ is typically achieved in
5--20~iterations.  Larger corrections are needed when composition
jumps are large (e.g.\ first timestep with a new core mass) or when
the rotation rate changes significantly.  The damping factor
$\alpha = 0.4$ sacrifices asymptotic quadratic convergence rate for
robustness: the standard $\alpha = 1$ (undamped) Newton method can
overshoot when the initial guess is far from the solution, particularly
for non-isothermal iron cores where the EOS curvature is strong.

\subsection{Unit System}

All internal computations use CGS units:
\begin{center}
\begin{tabular}{ll}
  \hline
  Quantity & Unit \\
  \hline
  Pressure & dyn\,cm$^{-2}$ \\
  Radius & cm \\
  Density & g\,cm$^{-3}$ \\
  Temperature & K \\
  Mass & g \\
  Entropy & $k_{\rm B}/$baryon \\
  \hline
\end{tabular}
\end{center}
Conversion factors are applied at the EOS interface boundaries:
$\times\,10^{-10}$ for dyn\,cm$^{-2}\to$\,GPa (mantle),
$\times\,10^{-1}$ for dyn\,cm$^{-2}\to$\,Pa (iron core).

\begin{acknowledgments}
ORCHARD builds on the APPLE code \citep{Sur2024}.
\end{acknowledgments}

\bibliography{hydrostatic_methods}
\bibliographystyle{aasjournal}

\end{document}
